\documentclass[11pt,a4paper]{article}
\usepackage[utf8]{inputenc}
\usepackage[T1]{fontenc} % ESSENCIAL: Resolve a quebra de caracteres no ATS
\usepackage{lmodern} % ESSENCIAL: Garante fonte mapeável e selecionável no PDF
\usepackage[margin=0.8in]{geometry}
\usepackage{titlesec}
\usepackage{enumitem}
\usepackage[hidelinks]{hyperref} % hidelinks remove a borda ao redor dos links

\titleformat{\section}{\Large\bfseries}{}{0em}{}[\titlerule]
\titlespacing{\section}{0pt}{1.5ex}{1ex}
\setlength{\parindent}{0pt}
\setlist[itemize]{leftmargin=*, nosep, itemsep=0.2em}

\begin{document}
\pagestyle{empty}

{\centering
\textbf{\huge Marcio Gabriel Schinor Mazega} \\ \vspace{0.2cm}
\textit{\large Engenheiro de Dados | Data Engineer} \\ \vspace{0.1cm}
Tupã, SP -- Brasil | (14) 99812-6535 | \href{mailto:marciomazegaa@gmail.com}{marciomazegaa@gmail.com} \\
\href{https://linkedin.com/in/marcio-mazega}{linkedin.com/in/marcio-mazega} | \href{https://github.com/marciomazega}{github.com/marciomazega} \\
}

\vspace{0.4cm}
\section*{RESUMO PROFISSIONAL}
Engenheiro de Dados com sólida experiência no desenho, orquestração e otimização de pipelines de ponta a ponta. Proficiente no desenvolvimento de rotinas de ETL/ELT escaláveis utilizando Python e na construção de arquiteturas de dados resilientes em PostgreSQL e MySQL. Foco em integrar múltiplas fontes (Web Scraping, APIs RESTful) para alimentar Data Lakes e Data Warehouses consistentes, aplicando rigor em governança e normalização relacional para viabilizar inteligência analítica e painéis de BI corporativo.

\vspace{0.2cm}
\section*{HABILIDADES TÉCNICAS}
\begin{itemize}
    \item \textbf{Engenharia de Dados:} ETL / Pipelines de Dados, Web Scraping, Integração de APIs, Data Quality.
    \item \textbf{Bancos de Dados e Infra:} PostgreSQL, MySQL, Docker, Linux (Ubuntu/WSL2).
    \item \textbf{Linguagens e Ferramentas:} Python (Pandas, Scikit-learn, FastAPI), SQL, Git / GitHub.
    \item \textbf{Arquitetura e Visualização:} Modelagem Relacional, APIs RESTful, Power BI, Dashboards.
\end{itemize}

\vspace{0.2cm}
\section*{EXPERIÊNCIA PROFISSIONAL}
\textbf{Engenheiro de Dados e Desenvolvedor Backend} \hfill \textit{Fev 2026 -- Jun 2026} \\
\textit{Gattaz Health} \\
\begin{itemize}
    \item Desenvolvi e orquestrei pipelines ETL ponta a ponta em Python, consolidando 100 registros/semana com alta eficiência.
    \item Automatizei a extração de dados via Web Scraping e estruturou o acesso via APIs RESTful (FastAPI), reduzindo o tempo de coleta manual em 50\%.
    \item Liderei o redesign físico e lógico do schema corporativo em PostgreSQL, otimizando as consultas para os painéis de BI e isolando dados clínicos sensíveis.
\end{itemize}

\vspace{0.2cm}
\textbf{Monitor de Iniciação Científica Júnior (ICJr)} \hfill \textit{Jan 2025 -- Mai 2026} \\
\textit{OBA - Olimpíada Brasileira de Astronomia e Astronáutica} \\
\begin{itemize}
    \item Orientei e supervisionei estudantes no desenvolvimento de projetos e correção de relatórios científicos (OBA/OBAFOG).
\end{itemize}

\vspace{0.2cm}
\textbf{Pesquisador de Dados AgTech (Bolsista FAPESP)} \hfill \textit{Ago 2025 -- Fev 2026} \\
\textit{Centro Avançado de Pesquisa Smart B100} \\
\begin{itemize}
    \item Aprendi sobre rigor analítico para organizar e classificar soluções tecnológicas, viabilizando métricas de competitividade para o setor.
\end{itemize}

\vspace{0.2cm}
\section*{PROJETOS DE DESTAQUE}
\textbf{PRF Analytics Full Cycle -- Pipeline ETL Histórico} | \href{https://github.com/marciomazega}{[Ver no GitHub]}
\begin{itemize}
    \item \textbf{Stack:} Python, Pandas, PostgreSQL, Web Scraping.
    \item \textbf{Descrição:} Pipeline completo (extração, tratamento, padronização) processando [X] anos de dados históricos da Polícia Rodoviária Federal, culminando em uma carga otimizada no PostgreSQL pronta para visualização.
\end{itemize}

\textbf{Buscador Inteligente de Artigos Científicos} | \href{https://github.com/marciomazega}{[Ver no GitHub]}
\begin{itemize}
    \item \textbf{Stack:} Python, NLP, FastAPI, Web Scraping.
    \item \textbf{Descrição:} Arquitetura automatizada para extração e filtragem de [X]+ publicações acadêmicas, conectando motores de scraping com uma API eficiente para otimizar fluxos de revisão bibliográfica.
\end{itemize}

\textbf{Central Agro -- Infraestrutura de Dados Geoespaciais}
\begin{itemize}
    \item \textbf{Stack:} Python, APIs Geoespaciais (Google Earth Engine).
    \item \textbf{Descrição:} Orquestração e transformação de camadas de dados climáticos e imagens de satélite em matrizes tabulares de [X] linhas, preparadas para ingestão em modelos de Machine Learning.
\end{itemize}

\vspace{0.2cm}
\section*{FORMAÇÃO ACADÊMICA E CERTIFICAÇÕES}
\textbf{Graduação em Tecnologia em Sistemas Inteligentes (TSI)} \hfill \textit{Jan 2025 -- Dez 2027} \\
\textit{FATEC Pompéia "Shunji Nishimura"} \\
\vspace{0.1cm}
\textbf{Certificações em Destaque:}
\begin{itemize}
    \item Introdução ao Machine Learning -- MBA USP/ESALQ (2026)
    \item Banco de Dados Relacional -- EncData
\end{itemize}

\vspace{0.2cm}
\section*{PREMIAÇÕES E PUBLICAÇÕES}
\begin{itemize}
    \item \textbf{Premiação:} 3º Lugar -- DSIN Coder Challenge (Desafio estadual de lógica e programação, 2025).
    \item \textbf{Publicações:} Coautor de 2 capítulos de livros acadêmicos focados em inovação no Agronegócio e Sociologia.
\end{itemize}

\end{document}